\documentclass{article}

% typesetting and text macros
\usepackage[letterpaper, textheight=9.0in, textwidth=5.2in]{geometry}
\usepackage{setspace}
\setstretch{1.08}
\usepackage{fancyhdr}
\setlength{\topmargin}{0.0in}
\setlength{\headheight}{3.0ex}\addtolength{\topmargin}{-0.5\headheight} % make room for headers
\setlength{\headsep}{3.0ex}\addtolength{\topmargin}{-0.5\headsep}
\pagestyle{fancy}
%\renewcommand{\headrulewidth}{0.0pt}
\fancyhf{}
\fancyhead[R]{\sffamily\thepage}
\sloppy\sloppypar\raggedbottom\frenchspacing

\fancyhead[L]{\sffamily Hogg / Self-calibration of pulsar timing arrays}
\title{\bfseries
Self-calibration of a pulsar timing array, in the context of a realistic gravitational-radiation sky}
\author{David W. Hogg}
\date{April 2026}

\begin{document}

\maketitle

\paragraph{Introduction:}
Right now, the way the pulsar-timing community looks for gravitational wave (GW) sources is by looking solely at the ``Earth term.''
That is, gravitational waves are found by looking for the correlations in timing residuals across pulsars from the fact that the Earth is being displaced by the GWs, relative to all of the pulsar clocks.
The pulsars are also being displaced of course, but those displacements are treated as just contributing some ``timing noise'' to the system.
When the dominant GW signal at the Earth is stochastic---which, for our purposes, means having significant contributions from large numbers of sources on the sky---the GW signal appears in the data only as correlations in residuals.
Things would be different if a single GW source dominated the timing signals, but we now know pretty confidently that this is not the case.

A (possibly irrelevant) historical fact about astronomy is that stochastic backgrounds have almost never been discovered prior to the discovery of discrete, identifiable sources.
For example, HOGG.

My position---controversial---is that \emph{if} we could phase up a pulsar timing array, the sensitivity of the array to GW sources would increase by an enormous factor, possibly $>10^7$.
Phasing up the array involves inferring the positions of the pulsars to a small fraction of a wavelenth, which is a positional precision of much less than a light year, for pulsars that are typically $>10^4$\,ly [HOGG?] away.
In principle the timing data themselves contain a lot of information about the pulsar positions; can we self-calibrate on that information?
Most approaches to this problem assume that there are very few sources, in part because the community is wedded to Bayesian approaches that are computationally intractable in any other situation.
I remind the reader that we know, at present, that there are \emph{not} very few sources dominating the timing-residual signals we have.

\paragraph{Problem setup:}

\paragraph{Correlations:}

\end{document}
